\section{Evaluation Setup: Supplementary Details}
\label{sec:appendix-setup}

This section consolidates setup details referenced from
\S\ref{sec:eval-method}, including synthetic workload
construction, operating-point selection, and scale and
microbenchmark details.

\subsection{Synthetic Workload Construction}
\label{sec:appendix-workloads}

The synthetic evaluation uses four workload families spanning
concentrated, diffuse, and heavy-tail noise workloads.
All workloads are generated over the same shared segment space;
drift is controlled independently by the side-imbalance
parameter~$q$.

\paragraph{Bimodal and diffuse}
\texttt{bimodal} concentrates all differences in $B$ dirty
segments, with the remaining segments clean.
\texttt{diffuse} spreads differences approximately uniformly
across a large number of segments.

\paragraph{Gray-zone heavy-tail workloads}
\texttt{heavy\_share} and \texttt{mix} place a large difference
count $d_{\text{big}}$ on a small set of segments while many
background segments each receive a smaller count
$d_{\text{small}}$.
For \texttt{mix}, difficulty is controlled by two axes:
noise-to-signal ratio $\texttt{S/B}$ and separability
$\texttt{ratio}{=}d_{\text{big}}/d_{\text{small}}$.
\texttt{heavy\_share} is the high-multiplier special case used
to probe behavior near the ranking-collapse boundary.

\paragraph{Recall target}
For mixed workloads, $\texttt{recall\_big}$ denotes the
fraction of segments carrying $d_{\text{big}}$ that are
covered by the candidate set.
This metric isolates whether the method recovers the
operationally important heavy segments when background noise is
present.

\subsection{Metrics and Operating-Point Protocol}
\label{sec:appendix-metrics}

At each stress point, every method sweeps its own internal
control knob:
SCOUT sweeps truncation ratio~\texttt{top\_r},
B2 sweeps top-$K$,
and B3 sweeps drill-down aggressiveness.
The operating point is the minimum candidate fraction
satisfying the recall requirement
$\mathtt{median(recall\_big)}\ge 0.95$.
If no setting satisfies the recall floor within the 2\%
candidate-budget cap, the point is recorded as
\textbf{N/A}.

\subsection{Scale and Microbenchmark Details}
\label{sec:appendix-scale}

Synthetic experiments sweep total segment count from
$10^4$ to $10^7$.
The YCSB validation replays up to ${\sim}2{\times}10^9$
records across $P{=}100{,}000$ range-partitioned segments,
producing nearly $50\times$ hot-to-cold segment-size skew
under the Zipfian workload.

\begin{table}[!thbp]
\centering
\small
\begin{tabular}{p{0.28\linewidth}p{0.62\linewidth}}
\toprule
Cost item & Reference anchor point \\
\midrule
\makecell[tl]{\textbf{Exchange} \\ (one-sided digest)}
  & $4VLm_0$ bytes; $V{=}2$: 512\,KiB, $V{=}3$: 768\,KiB \\
\makecell[tl]{\textbf{Decode} \\ (stats + score)}
  & $\sim$0.33\,ms/view; $\sim$1\,ms at $V{=}3$; independent of $P$ \\
\makecell[tl]{\textbf{Build} \\ (per-key update)}
  & $\sim$2.9\,$\mu$s/key; includes window lookups and $k$ hash updates \\
\bottomrule
\end{tabular}
\caption{Reference anchor points for probe overhead ($L{=}2048$, $m_0{=}32$, $k{=}2$).
Decode and Build are measured from a C++ microbenchmark.}
\label{tab:cost-anchor}
\end{table}

At this scale, the probe's digest-side overhead remains small:
one-sided digest size stays below 1\,MiB and decode time stays
below 1\,ms; reference anchor points for exchange size, decode
latency, and build cost are summarized in
Table~\ref{tab:cost-anchor}.
In a single-threaded microbenchmark on a Xeon Platinum 8468,
each key update takes approximately $2.9\,\mu$s; building the
probe over $10^9$ keys therefore completes in minutes, with
the dominant wall-clock cost still determined by local
sequential-scan I/O rather than probe computation.
